% !TeX program = xelatex
% Compile with XeLaTeX in Overleaf.


\documentclass[a4paper]{article}
\usepackage{xeCJK}
\setCJKmainfont{FandolSong-Regular}
\usepackage{needspace}

% \usepackage[default]{fontsetup}

\usepackage{fancyhdr}
\usepackage{extramarks}
\usepackage{amsmath}
\usepackage{amsthm}
\usepackage{amsfonts}
\usepackage{tikz}
\usepackage[plain]{algorithm}
\usepackage{algpseudocode}
\usepackage{enumerate}
\usepackage{tikz}

\usetikzlibrary{automata,positioning}

%
% Basic Document Settings
%  

\topmargin=-0.2in
\evensidemargin=0in
\oddsidemargin=0in
\textwidth=6.5in
\textheight=9.5in
\headsep=0.25in

\linespread{1.1}

\pagestyle{fancy}
\lhead{}
\chead{\hmwkClass : \hmwkTitle}
\rhead{\firstxmark}
\lfoot{\lastxmark}
\cfoot{\thepage}

\renewcommand\headrulewidth{0.4pt}
\renewcommand\footrulewidth{0.4pt}

\setlength\parindent{0pt}

%
% Create Problem Sections
%

\newcommand{\enterProblemHeader}[1]{
    \nobreak\extramarks{}{Problem \arabic{#1} continued on next page\ldots}\nobreak{}
    \nobreak\extramarks{Problem \arabic{#1} (continued)}{Problem \arabic{#1} continued on next page\ldots}\nobreak{}
}

\newcommand{\exitProblemHeader}[1]{
    \nobreak\extramarks{Problem \arabic{#1} (continued)}{Problem \arabic{#1} continued on next page\ldots}\nobreak{}
    \stepcounter{#1}
    \nobreak\extramarks{Problem \arabic{#1}}{}\nobreak{}
}

\newcommand*\circled[1]{\tikz[baseline=(char.base)]{
		\node[shape=circle,draw,inner sep=2pt] (char) {#1};}}


\setcounter{secnumdepth}{0}
\newcounter{partCounter}
\newcounter{homeworkProblemCounter}
\setcounter{homeworkProblemCounter}{1}
\nobreak\extramarks{Problem \arabic{homeworkProblemCounter}}{}\nobreak{}

%
% Homework Problem Environment
%
% This environment takes an optional argument. When given, it will adjust the
% problem counter. This is useful for when the problems given for your
% assignment aren't sequential. See the last 3 problems of this template for an
% example.
%

\newenvironment{homeworkProblem}[1][-1]{
    \ifnum#1>0
        \setcounter{homeworkProblemCounter}{#1}
    \fi
    \section{Problem \arabic{homeworkProblemCounter}}
    \setcounter{partCounter}{1}
    \enterProblemHeader{homeworkProblemCounter}
}{
    \exitProblemHeader{homeworkProblemCounter}
}

%
% Homework Details
%   - Title
%   - Class
%   - Due date
%   - Name
%   - Student ID

\newcommand{\hmwkTitle}{Homework\ \#02}
\newcommand{\hmwkClass}{Probability \& Statistics for EECS}
\newcommand{\hmwkDueDate}{2024-10-6}
\newcommand{\hmwkAuthorName}{杨涵}
\newcommand{\hmwkAuthorID}{2025531064}


%
% Title Page
%

\title{
    \vspace{2in}
    \textmd{\textbf{\hmwkClass:\\  \hmwkTitle}}\\
    \normalsize\vspace{0.1in}\small{Due\ on\ \hmwkDueDate\ at 18:00}\\
	\vspace{4in}
}

\author{
	Name: \textbf{\hmwkAuthorName} \\
	Student ID: \hmwkAuthorID}
\date{}

\renewcommand{\part}[1]{\textbf{\large Part \Alph{partCounter}}\stepcounter{partCounter}\\}

%
% Various Helper Commands
%

% Useful for algorithms
\newcommand{\alg}[1]{\textsc{\bfseries \footnotesize #1}}
% For derivatives
\newcommand{\deriv}[1]{\frac{\mathrm{d}}{\mathrm{d}x} (#1)}
% For partial derivatives
\newcommand{\pderiv}[2]{\frac{\partial}{\partial #1} (#2)}
% Integral dx
\newcommand{\dx}{\mathrm{d}x}
% Alias for the Solution section header
\newcommand{\solution}{\par\addvspace{\baselineskip}\Needspace{4\baselineskip}\noindent\textbf{\large Solution}\par\nobreak\smallskip}
% Probability commands: Expectation, Variance, Covariance, Bias
\newcommand{\E}{\mathrm{E}}
\newcommand{\Var}{\mathrm{Var}}
\newcommand{\Cov}{\mathrm{Cov}}
\newcommand{\Bias}{\mathrm{Bias}}

\begin{document}


% \maketitle
% \thispagestyle{empty}
% \pagebreak

\date{
Due on Oct. 7, 2026, 23:59 UTC+8}
\title{SI 140A-01  Probability \& Statistics for EECS, Fall 2026 \\
Homework 2}
\maketitle
Read all the instructions below carefully before you start working on the assignment, and before you make a submission.
\begin{itemize}
    \item You are required to write down all the major steps towards making your conclusions; otherwise you may obtain limited points of the problem.
    \item Write your homework in English; otherwise you will get no points of this homework.
    \item Any form of plagiarism will lead to $0$ point of this homework. 
\end{itemize}
\newpage

\begin{homeworkProblem}[1]
Show that for any events $A$ and $B$,
    \[
    P(A) + P(B) - 1 \le P(A \cap B) \le P(A \cup B) \le P(A) + P(B).
    \]
    For each of these three inequalities, give a simple criterion for when the inequality is actually an equality (e.g., give a simple condition such that $P(A \cap B) = P(A \cup B)$ if and only if the condition holds).

\solution
According to the inclusion--exclusion formula,
\[
P(A\cup B)=P(A)+P(B)-P(A\cap B).
\]
\begin{enumerate}[(i)]
\item Since $P(A\cup B)\leq 1$,
\[
P(A)+P(B)-P(A\cap B)\leq 1
\quad\Longrightarrow\quad
P(A)+P(B)-1\leq P(A\cap B).
\]
Equality holds if and only if $P(A\cup B)=1$.

\item Since $A\cap B\subseteq A\cup B$,
\[
P(A\cap B)\leq P(A\cup B).
\]
Equality holds if and only if $P(A\triangle B)=0$.

\item Since $P(A\cap B)\geq 0$,
\[
P(A)+P(B)-P(A\cup B)\geq 0
\quad\Longrightarrow\quad
P(A\cup B)\leq P(A)+P(B).
\]
Equality holds if and only if $P(A\cap B)=0$.
\end{enumerate}

\end{homeworkProblem}

\newpage
\begin{homeworkProblem}[2]
Let $A$ and $B$ be events. The \emph{symmetric difference} $A \triangle B$ is defined to be the set of all elements that are in $A$ or $B$ but not both. In logic and engineering, this event is also called the $\text{XOR}$ (\emph{exclusive or}) of $A$ and $B$. Show that
    \[
    P(A \triangle B) = P(A) + P(B) - 2P(A \cap B),
    \]
    directly using the axioms of probability.

\solution
By definition,
\[
A\triangle B=(A\setminus B)\cup(B\setminus A).
\]
The events $A\setminus B$ and $B\setminus A$ are disjoint, so
\[
P(A\triangle B)=P(A\setminus B)+P(B\setminus A).
\]
As
\[
A=(A\setminus B)\cup(A\cap B),
\qquad B=(B\setminus A)\cup(A\cap B),
\]
where each union is disjoint, we have
\[
P(A\setminus B)=P(A)-P(A\cap B),
\qquad P(B\setminus A)=P(B)-P(A\cap B).
\]
Therefore,
\begin{align*}
P(A\triangle B)
&=P(A)-P(A\cap B)+P(B)-P(A\cap B)\\
&=P(A)+P(B)-2P(A\cap B).
\end{align*}

\end{homeworkProblem}



%\newpage
%\begin{homeworkProblem}[3]
%Given $n \geq 2$ different numbers $\left(a_1, a_2, \ldots, a_n\right)$, a bootstrap sample is a sequence $\left(x_1, x_2, \ldots, x_n\right)$ formed from the $a_j$ 's by sampling with replacement with equal probabilities.  For example, if $n=2$ and $\left(a_1, a_2\right)=(3,1)$, then the possible bootstrap samples are $(3,3),(3,1),(1,3)$, and $(1,1)$.\\
%(a) How many possible bootstrap samples are there for $\left(a_1, \ldots, a_n\right)$ ?\\
%(b) How many possible bootstrap samples are there for $\left(a_1, \ldots, a_n\right)$, if order does not matter.\\ %(in the sense that it only matters how many times each $a_j$ was chosen, not the order in which they were chosen)?\\
%(c) One random bootstrap sample is chosen (by sampling from $a_1, \ldots, a_n$ with replacement, as described above). Show that not all unordered bootstrap samples (in the sense of (b)) are equally likely. Find an unordered bootstrap sample $\mathbf{b}_1$ that is as likely as possible, and an unordered bootstrap sample $\mathbf{b}_2$ that is as unlikely as possible. Let $p_1$ be the probability of getting $\mathbf{b}_1$ and $p_2$ be the probability of getting $\mathbf{b}_2$ (so $p_i$ is the probability of getting the specific unordered bootstrap sample $\mathbf{b}_i$ ). What is $p_1 / p_2$ ? What is the ratio of the probability of getting an unordered bootstrap sample whose probability is $p_1$ to the probability of getting an unordered sample whose probability is $p_2$ ?

%\end{homeworkProblem}

% \newpage
% \begin{homeworkProblem}[4]
% \textbf{(Geometric Probability)} You get a stick and break it randomly into three pieces. What is the probability that you can make a triangle using such three pieces?
% \end{homeworkProblem}


\newpage
\begin{homeworkProblem}[3]
Events $A$ and $B$ are \emph{independent} if $P(A \cap B) = P(A)P(B)$.
    \begin{enumerate}
        \item Give an example of independent events $A$ and $B$ in a finite sample space $S$ (with neither equal to $\emptyset$ or $S$).
        \item Consider the experiment of picking a random point in the rectangle
        \[
        R = \{(x, y) : 0 < x < 1, 0 < y < 1\},
        \]
        where the probability of the point being in any particular region contained within $R$ is the area of that region. Let $A_1$ and $B_1$ be rectangles contained within $R$, with areas not equal to 0 or 1. Let $A$ be the event that the random point is in $A_1$, and $B$ be the event that the random point is in $B_1$. Give a geometric description of when it is true that $A$ and $B$ are independent. Also, give an example where they are independent and another example where they are not independent.
        \item Show that if $A$ and $B$ are independent, then
        \[
        P(A \cup B) = P(A) + P(B) - P(A)P(B) = 1 - P(A^c)P(B^c).
        \]
    \end{enumerate}


\solution
\begin{enumerate}[(i)]
\item Assume $S=\{1,2,3,4\}$, $A=\{1,2\}$, and $B=\{2,4\}$. Then
\[
P(A)=\frac12,\qquad P(B)=\frac12,\qquad A\cap B=\{2\},
\]
and
\[
P(A\cap B)=\frac14=P(A)P(B).
\]
Thus $A$ and $B$ are independent.

\item Let $S_{A_1}$, $S_{B_1}$, and $S_{A_1\cap B_1}$ denote the corresponding areas. Then
\[
P(A)=S_{A_1},\qquad P(B)=S_{B_1},\qquad
P(A\cap B)=S_{A_1\cap B_1}.
\]
The events $A$ and $B$ are independent when
\[
S_{A_1}S_{B_1}=S_{A_1\cap B_1}.
\]
Geometrically, the overlap of $A_1$ and $B_1$ should have an area equal to the product of their individual areas.

\medskip
\Needspace{9\baselineskip}
For an independent example, let
\[
A_1=\{(x,y):0<x<\tfrac12,\ 0<y<1\},\qquad
B_1=\{(x,y):0<x<1,\ 0<y<\tfrac12\}.
\]
Then
\[
P(A)=\frac12,\qquad P(B)=\frac12,\qquad
P(A\cap B)=\frac14=P(A)P(B).
\]

\medskip
\Needspace{8\baselineskip}
For a non-independent example, let
\[
A_1=B_1=\{(x,y):0<x<\tfrac12,\ 0<y<1\}.
\]
Then
\[
P(A)=P(B)=\frac12,\qquad
P(A\cap B)=\frac12\neq P(A)P(B).
\]

\item Since $A$ and $B$ are independent, $P(A\cap B)=P(A)P(B)$. Thus
\[
P(A\cup B)=P(A)+P(B)-P(A\cap B)
=P(A)+P(B)-P(A)P(B).
\]
Also,
\[
P(A^c)P(B^c)=P((A\cup B)^c)=1-P(A\cup B),
\]
so
\[
P(A\cup B)=1-P(A^c)P(B^c).
\]
\end{enumerate}

\end{homeworkProblem}

% \newpage
% \begin{homeworkProblem}[4]
% \textbf{(Geometric Probability)} You get a stick and break it randomly into three pieces. What is the probability that you can make a triangle using such three pieces?
% \end{homeworkProblem}

\newpage
\begin{homeworkProblem}[4]
Consider four nonstandard dice (the \emph{Efron dice}), whose sides are labeled as follows (the 6 sides on each die are equally likely).
    \begin{itemize}
        \item[] A: $4, 4, 4, 4, 0, 0$
        \item[] B: $3, 3, 3, 3, 3, 3$
        \item[] C: $6, 6, 2, 2, 2, 2$
        \item[] D: $5, 5, 5, 1, 1, 1$
    \end{itemize}
    These four dice are each rolled once. Let $A$ be the result for die $A$, $B$ be the result for die $B$, etc.
    \begin{enumerate}
        \item Find $P(A > B)$, $P(B > C)$, $P(C > D)$, and $P(D > A)$.
        \item Is the event $A > B$ independent of the event $B > C$? Is the event $B > C$ independent of the event $C > D$? Explain.
    \end{enumerate}

\solution
\begin{enumerate}[(i)]
\item The event $A>B$ occurs when $A=4$, so
\[
P(A>B)=\frac46=\frac23.
\]
The event $B>C$ occurs when $C=2$, so
\[
P(B>C)=\frac46=\frac23.
\]
The event $C>D$ occurs when $C=6$, or when $C=2$ and $D=1$. These cases have probabilities
\[
p_1=\frac26=\frac13,\qquad
p_2=\frac23\cdot\frac12=\frac13.
\]
Hence
\[
P(C>D)=\frac13+\frac13=\frac23.
\]
The event $D>A$ occurs when $D=5$, or when $D=1$ and $A=0$. These cases have probabilities
\[
p_1=\frac12,\qquad p_2=\frac12\cdot\frac13=\frac16.
\]
Hence
\[
P(D>A)=\frac12+\frac16=\frac23.
\]

\item We have $P(A>B)=P(B>C)=\frac23$. The event
$\{A>B\}\cap\{B>C\}$ occurs when $A=4$, $B=3$, and $C=2$. Thus
\[
P(\{A>B\}\cap\{B>C\})
=\frac23\cdot1\cdot\frac23=\frac49
=P(A>B)P(B>C).
\]
Therefore, the events $A>B$ and $B>C$ are independent.

\medskip
\Needspace{8\baselineskip}
We also have $P(B>C)=P(C>D)=\frac23$. The event
$\{B>C\}\cap\{C>D\}$ occurs when $B=3$, $C=2$, and $D=1$. Thus
\[
P(\{B>C\}\cap\{C>D\})
=1\cdot\frac23\cdot\frac12=\frac13
\neq P(B>C)P(C>D).
\]
Therefore, the events $B>C$ and $C>D$ are not independent, because $B>C$ implies $C=2$, which changes the probability that $C>D$.
\end{enumerate}

\end{homeworkProblem}

 




\end{document}
